\scanpage{7}
\begin{center}( 2 )\end{center}

(ii) En caractéristique $p>0$ (pour les corps résiduels) : si $B$ est quasi-séparable, $L_b$ est relativement parfait \textsuperscript{(2)} sur $L_a$ et
\[
(\mathrm S)\quad \mathfrak n=(\mathfrak n\cap A[B^p])B+\mathfrak n^2,
\]
la réciproque étant vraie si l'on suppose en outre $D_A(B)$ de type fini. Si $B$ est noethérien, (S) se simplifie encore.

\textit{Globalement :}

(i) Soit $A$ un anneau noethérien quasi-séparable sur un corps $K$ :

--- en caractéristique $0$, $A$ est composé direct d'une famille finie d'extensions algébriques de $K$;

--- en caractéristique $p>0$, si pour tout idéal maximal de $A$ le corps résiduel est une extension séparable (resp. de type fini) de $K$, $A$ est composé direct d'une famille finie d'extensions séparables et quasi-séparables (resp. algébriques séparables) de $K$.

(ii) Soit $A$ une algèbre quasi-séparable sur un corps $K$ :

--- en caractéristique $0$, tout idéal premier de $A$ est maximal et tous les corps résiduels des localisés de $A$ sont des extensions algébriques de $K$;

--- en caractéristique $p>0$, si pour tout idéal premier $\unsure{\mathfrak q|\mathfrak g}$ de $A$ le corps $A_{\unsure{\mathfrak q}}/\unsure{\mathfrak q}A_{\unsure{\mathfrak q}}$ est de type fini sur $K$, alors tout idéal premier de $A$ est maximal et tous les corps résiduels sont des extensions algébriques séparables de $K$;

--- Sous l'une ou l'autre de ces hypothèses, le quotient de $A$ par sa nilracine est un anneau absolument plat \textsuperscript{(3)}.

2. \textsc{Différentielles et anneaux locaux réguliers.} --- On se place dans le cadre local, en supposant même $A$ et $B$ intègres, de corps des fractions $K_a$ et $K_b$, et les cardinaux $\{L_b:L_a\}$, $\{K_b:K_a\}$, $\operatorname{rg}\mathfrak n$ finis.

\textsc{Théorème.} --- \textit{Si $D_A(B)$ est de type fini (en particulier si $B$ est de type fini sur $A$) il y a équivalence entre :}

(i) \textit{$D_A(B)$ est libre;}

(ii) \textit{a. Dans le cas où $L_b$ est séparable sur $L_a$ :}
\[
\dim_{L_b}\frac{\mathfrak n}{\mathfrak n^2+\mathfrak mB}=\{K_b:K_a\}-\{L_b:L_a\};
\]

\textit{b. Sinon, $p>0$ désignant la caractéristique de $L_b$ et $L_a$ :}
\[
\dim_{L_b}\frac{\mathfrak n}{\mathfrak n^2+(\mathfrak n\cap A[B^p])B}=\{K_b:K_a\}-\{L_b:L_a\}.
\]

Notons que l'hypothèse « $D_A(B)$ de type fini » entraîne que les cardinaux $\{K_b:K_a\}$ et $\{L_b:L_a\}$ sont finis. De plus, l'égalité (ii) $a$ s'écrit :

--- Si $A$ est un corps : $\operatorname{rang}\mathfrak n=\{K_b:A\}-\{L_b:A\}$,

--- Si $B$ est non ramifié sur $A$ : $\{K_b\mid K_a\}=\{L_b\mid L_a\}$.

\mspage{6}
